% ---------------------------------------------------------------------
%  CHƯƠNG 2 — PHÂN TÍCH, THIẾT KẾ VÀ TRIỂN KHAI
% ---------------------------------------------------------------------
\chapter{Phân tích, thiết kế và triển khai}
\label{ch:thietke}

\begin{muctieu}
Chương này phân tích vai trò và yêu cầu, mô tả kiến trúc, dữ liệu, quy
trình sinh câu hỏi và tổ chức thi, rồi trình bày cài đặt và thử nghiệm.
\end{muctieu}

\section{Phân tích bài toán và yêu cầu}

Hệ thống có ba vai trò chính: quản trị viên quản lý tài khoản và dữ liệu
toàn hệ thống; giảng viên quản lý câu hỏi, tài liệu, đề thi, lớp học và
chấm bài; sinh viên tham gia lớp/phòng thi và xem lịch sử dự thi. Các yêu
cầu chức năng trọng tâm được tóm tắt trong Bảng~\ref{tab:yeucau}.

\begin{table}[H]
\centering
\caption{Nhóm yêu cầu chức năng của hệ thống.}
\label{tab:yeucau}
\begin{tabular*}{\textwidth}{@{\extracolsep{\fill}}p{0.12\textwidth}p{0.31\textwidth}p{0.49\textwidth}}
\toprule
\rowcolor{gray100}\bfseries\color{daunavy}Mã & \bfseries\color{daunavy}Chức năng & \bfseries\color{daunavy}Mô tả \\
\midrule
FR1 & Tài khoản và phân quyền & Đăng nhập, hồ sơ người dùng và giới hạn chức năng theo vai trò. \\
FR2 & Ngân hàng câu hỏi & Tạo, sửa, xoá, lọc và phân loại câu hỏi theo môn, loại và độ khó. \\
FR3 & Hỗ trợ tạo đề & Đọc tài liệu, đề xuất câu hỏi, chọn phân bố, lưu đề và tạo mã đề. \\
FR4 & Duyệt đề & Chuyển đề qua các trạng thái nháp, chờ duyệt, đã duyệt và phát hành. \\
FR5 & Thi trực tuyến & Tạo phòng/mã tham gia, lưu câu trả lời, tiếp tục lượt thi và tự chấm trắc nghiệm. \\
FR6 & Lớp học & Tạo lớp, tham gia bằng mã, quản lý trạng thái buổi học và thành viên. \\
FR7 & Chấm và thống kê & Giảng viên chấm bài, ghi nhận phản hồi và xem số liệu theo quyền. \\
\bottomrule
\end{tabular*}
\end{table}

Các yêu cầu phi chức năng gồm bảo vệ API bằng xác thực và phân quyền, lưu
dữ liệu bền vững qua PostgreSQL, kiểm tra đầu vào ở các luồng quan trọng,
và hỗ trợ chạy các thành phần chính bằng Docker Compose. Báo cáo chưa xác
nhận ngưỡng thời gian phản hồi hay số người dùng đồng thời; các tiêu chí đó
cần được đo riêng trước triển khai.

\subsection{Tác nhân và kịch bản sử dụng}
Các tác nhân tương tác với hệ thống qua giao diện web. Quản trị viên tạo và
quản lý tài khoản, phân quyền và xem số liệu toàn hệ thống. Giảng viên nạp
học liệu, xem các câu hỏi được đề xuất, hiệu chỉnh câu hỏi, lưu thành ngân
hàng, tạo/duyệt đề, mở phòng thi, giám sát, chấm bài và xem phân tích.
Sinh viên đăng nhập hoặc tham gia lớp/phòng thi theo quyền, trả lời câu
hỏi, đánh dấu nội dung cần xem lại và nộp bài. Một lượt dự thi được gắn với
phòng, đề và tài khoản sinh viên.

Bảng~\ref{tab:usecase} mô tả một số kịch bản trọng yếu. Kịch bản thể hiện
hành vi mong muốn của sản phẩm; các bước được triển khai đến đâu được đánh
giá riêng trong Chương~\ref{ch:ketqua}.

\begin{table}[H]
\centering
\caption{Đặc tả các kịch bản sử dụng chính.}
\label{tab:usecase}
\begin{tabularx}{\textwidth}{>{\bfseries\color{daunavy}}p{0.16\textwidth} X X}
\toprule
\rowcolor{gray100}Kịch bản & Luồng chính & Kết quả / ngoại lệ cần xử lý \\
\midrule
UC1 -- Sinh câu hỏi & Giảng viên chọn tài liệu, môn học, số lượng, dạng câu và phân bố; gửi phân tích; xem ứng viên, bằng chứng và hiệu chỉnh. & Ứng viên có nguồn để đối chiếu; từ chối tài liệu rỗng, lỗi đọc hoặc đầu ra không hợp lệ. \\
UC2 -- Tạo/phát hành đề & Chọn câu hỏi; khai báo thời lượng, thang điểm và metadata; lưu nháp; gửi duyệt; phê duyệt và phát hành. & Đề lưu kèm liên kết câu hỏi; lỗi ma trận không đủ câu phải được thông báo. \\
UC3 -- Tổ chức thi & Giảng viên chọn đề, đặt thời gian, tạo mã; sinh viên nhập mã, làm bài, lưu và nộp. & Một sinh viên tiếp tục đúng lượt thi đã bắt đầu; bài nộp được tự chấm trắc nghiệm. \\
UC4 -- Chấm tự luận & Giảng viên chọn bài đã nộp, nhập điểm trong giới hạn và phản hồi, lưu kết quả. & Kết quả gắn với bài nộp; không cho điểm vượt tổng điểm đề. \\
UC5 -- Tham gia lớp & Giảng viên tạo lớp, chọn kiểu kết nối; sinh viên nhập mã; giảng viên mở/đóng buổi học. & Thành viên được lưu; lớp đã kết thúc hoặc đóng nhận thành viên thì từ chối lượt vào mới. \\
\bottomrule
\end{tabularx}
\end{table}

\begin{figure}[H]
\centering
\resizebox{0.9\textwidth}{!}{%
\begin{tikzpicture}[node distance=8mm and 12mm]
  \node[objbox] (teacher) {Giảng viên};
  \node[objbox, below=of teacher] (student) {Sinh viên};
  \node[objbox, above=of teacher] (admin) {Quản trị viên};
  \node[clsbox, right=of teacher] (bank) {Quản lý ngân hàng\\câu hỏi};
  \node[clsbox, right=of bank] (exam) {Tạo / duyệt\\đề thi};
  \node[clsbox, right=of exam] (room) {Thi và giám sát\\phòng thi};
  \node[clsbox, below=of exam] (classroom) {Lớp học\\trực tuyến};
  \node[clsbox, below=of bank] (grading) {Chấm bài\\và thống kê};
  \draw[flowarr] (teacher) -- (bank);
  \draw[flowarr] (teacher) -- (exam);
  \draw[flowarr] (teacher) -- (room);
  \draw[flowarr] (teacher) -- (grading);
  \draw[flowarr] (teacher) -- (classroom);
  \draw[flowarr] (student) -- (room);
  \draw[flowarr] (student) -- (classroom);
  \draw[flowarr] (student) -- (grading);
  \draw[flowarr] (admin) -- (bank);
  \draw[flowarr] (admin) -- (exam);
  \draw[flowarr] (admin) -- (room);
\end{tikzpicture}}
\caption{Quan hệ giữa các vai trò và nhóm chức năng ở mức use case.}
\label{fig:use-case}
\end{figure}
Hình~\ref{fig:use-case} chỉ thể hiện nhóm chức năng theo vai trò, không thay
thế bảng endpoint hoặc chính sách phân quyền chi tiết.

\subsection{Yêu cầu phi chức năng và tiêu chí nghiệm thu}
Yêu cầu phi chức năng được phân thành an toàn, toàn vẹn dữ liệu, khả dụng,
dễ vận hành và khả năng tiếp cận. Trạng thái đạt/đo được trình bày ở
Bảng~\ref{tab:nfr}; các ngưỡng chưa đo không được xem là đã thoả mãn.
\begin{table}[H]
\centering
\caption{Yêu cầu phi chức năng và mức xác minh hiện có.}
\label{tab:nfr}
\begin{tabularx}{\textwidth}{>{\bfseries\color{daunavy}}p{0.20\textwidth} X p{0.23\textwidth}}
\toprule
\rowcolor{gray100}Thuộc tính & Tiêu chí yêu cầu & Bằng chứng hiện có \\
\midrule
Phân quyền & Mỗi thao tác ghi dữ liệu chỉ được phép với vai trò và chủ sở hữu thích hợp. & Có middleware và test phân quyền; chưa có kiểm thử độc lập đầy đủ. \\
Toàn vẹn & Khóa ngoại, ràng buộc, giao dịch tạo đề/phòng và tính duy nhất. & Có schema, migration và transaction ở các luồng chính. \\
Khả dụng & Lỗi dịch vụ ngoài không được biến thành kết quả thành công giả. & Có fallback cho sinh nội dung; cần kiểm tra log và UX ở triển khai thật. \\
Hiệu năng & Giao diện và API đáp ứng thời gian chấp nhận được ở tải mục tiêu. & Chưa có ngưỡng hoặc đo tải. \\
Khả năng tiếp cận & Điều khiển bằng bàn phím, nhãn biểu mẫu, trạng thái và tương phản phù hợp. & Chưa thực hiện đánh giá WCAG có hệ thống~\cite{WCAG22}. \\
Quản lý dữ liệu & Hạn chế dữ liệu cá nhân/tài liệu gửi cho bên thứ ba; quy định lưu trữ và xoá. & Chưa có kiểm toán quyền riêng tư hoặc chính sách vận hành hoàn chỉnh. \\
\bottomrule
\end{tabularx}
\end{table}

\section{Thiết kế kiến trúc hệ thống}

Kiến trúc triển khai được trình bày ở Hình~\ref{fig:kientruc}. Frontend
React gửi yêu cầu HTTP đến REST API Express; backend áp dụng xác thực,
nghiệp vụ và truy vấn PostgreSQL. Luồng tài liệu đi qua bộ trích xuất văn
bản, sau đó có thể gọi dịch vụ mô hình ngôn ngữ ngoài hoặc bộ sinh dự phòng.
Các thư mục \code{ai-services/} và \code{ml-models/} hiện là khung tổ chức,
không phải các dịch vụ độc lập đang được gọi trong luồng runtime.

\begin{figure}[H]
\centering
\resizebox{\textwidth}{!}{%
\begin{tikzpicture}[node distance=12mm and 14mm]
  \node[objbox] (user) {Giảng viên /\\ sinh viên / admin};
  \node[clsbox, right=of user] (web) {Frontend\\ React + TypeScript};
  \node[clsbox, right=of web] (api) {REST API\\ Node.js + Express};
  \node[clsbox, right=of api] (db) {PostgreSQL\\ users, exams, attempts};
  \node[clsbox, below=of api] (doc) {Trích xuất văn bản\\ PDF / DOCX / OCR};
  \node[objbox, right=of doc] (ai) {OpenAI API tùy chọn\\ hoặc bộ sinh dự phòng};
  \draw[flowarr] (user) -- (web);
  \draw[flowarr] (web) -- (api);
  \draw[flowarr] (api) -- (db);
  \draw[flowarr] (api) -- (doc);
  \draw[flowarr] (doc) -- (ai);
  \draw[flowarr] (ai.north) |- (api.south);
\end{tikzpicture}}
\caption{Kiến trúc logic và luồng xử lý chính của AI ExamHub DAU.}
\label{fig:kientruc}
\end{figure}

\subsection{Phân rã thành phần}
Frontend đóng vai trò client hiển thị view theo vai trò và gọi các hàm
trong lớp dịch vụ API. Backend tập trung xử lý các quy tắc miền, xác thực
JWT, kiểm tra vai trò, thao tác cơ sở dữ liệu và tích hợp mô hình ngoài.
PostgreSQL giữ dữ liệu nghiệp vụ; Docker Compose nối backend, frontend và
cơ sở dữ liệu trong môi trường phát triển. Các module OCR/RAG/embedding
trong cây thư mục dự án là vị trí quy hoạch cho tương lai, không đồng nghĩa
các container AI này đã tồn tại trong cấu hình Compose hiện tại.

\subsection{Luồng xử lý phân tích tài liệu}
Luồng phân tích nhận \code{documentId}, môn học, số lượng, nội dung, ma
trận loại câu, phân bố độ khó và mức độ nhận thức. Backend từ chối nội
dung rỗng; nếu có nội dung, bộ phân tích tách các câu hỏi được định dạng
sẵn, đồng thời gọi bộ sinh từ tài liệu. Các ứng viên được hợp nhất, chuẩn
hoá và lọc theo ma trận. Những câu phù hợp được lưu vào ngân hàng câu hỏi
trước khi trả về giao diện. Mô hình dữ liệu do đó phân biệt việc phân tích
đề xuất với việc người dùng chọn câu để tạo đề hoàn chỉnh.

\subsection{Luồng thi trực tuyến}
Giảng viên tạo phòng từ một hoặc nhiều đề đã lưu. Mã phòng được tạo ngẫu
nhiên; sinh viên dùng mã để bắt đầu lượt thi. Backend ghi một lượt cho
cặp phòng--sinh viên, bộ câu hỏi được giao, đáp án, câu đánh dấu, vị trí
hiện tại, thời hạn và thời điểm nộp. Khi nộp, bài thi trắc nghiệm được
tự chấm; đề tự luận có thể yêu cầu chấm riêng. Giảng viên có thể xem tiến
độ và sự kiện chuyển tab được ghi lại. Sự kiện chuyển tab là tín hiệu hỗ
trợ giám sát, không đủ để tự kết luận có gian lận.

\section{Thiết kế dữ liệu}

Hệ thống không sử dụng tập dữ liệu huấn luyện phân lớp. Dữ liệu nghiệp vụ
gồm tài khoản, câu hỏi, đề thi, lượt dự thi, bài nộp, kết quả và lớp học.
Các bảng lõi được mô tả khái quát ở Bảng~\ref{tab:dulieu}; các trường chi
tiết được khai báo trong lược đồ PostgreSQL của dự án.

\begin{table}[H]
\centering
\caption{Các thực thể chính trong cơ sở dữ liệu.}
\label{tab:dulieu}
\begin{tabular*}{\textwidth}{@{\extracolsep{\fill}}>{\bfseries\color{daunavy}}p{0.19\textwidth}>{\color{daumaroon}}p{0.28\textwidth}p{0.43\textwidth}}
\toprule
\rowcolor{gray100}\bfseries\color{daunavy}Thực thể & \bfseries\color{daumaroon}Bảng & \bfseries\color{daunavy}Vai trò \\
\midrule
Tài khoản & users & Email, vai trò, hồ sơ và thông tin định danh. \\
Câu hỏi & questions & Nội dung, lựa chọn, đáp án, bằng chứng nguồn và rubric. \\
Đề thi & exams, exam\_questions & Metadata đề và liên kết có thứ tự tới câu hỏi. \\
Thi trực tuyến & exam\_rooms, exam\_attempts & Phòng thi, câu hỏi được giao, đáp án và thời điểm. \\
Bài nộp/kết quả & submissions, results & Trạng thái nộp, điểm số và phản hồi. \\
Lớp học & classroom\_rooms, classroom\_members & Buổi học, mã tham gia và thành viên. \\
\bottomrule
\end{tabular*}
\end{table}

\subsection{Sơ đồ thực thể--liên kết}
Các quan hệ lõi trong schema được tóm tắt tại Hình~\ref{fig:erd}. Người
dùng tạo câu hỏi/đề/phòng/lớp; đề liên kết đến câu hỏi qua bảng nối; lượt
thi thuộc về một sinh viên, một phòng và một đề; bài nộp liên kết với kết
quả. Bảng lớp học nối người dùng sinh viên với phòng học bằng khoá ghép.
\begin{figure}[H]
\centering
\resizebox{\textwidth}{!}{%
\begin{tikzpicture}[node distance=10mm and 13mm]
  \node[clsbox] (users) {users};
  \node[clsbox, right=of users] (questions) {questions};
  \node[clsbox, right=of questions] (exams) {exams};
  \node[objbox, below=of questions] (examq) {exam\_questions};
  \node[clsbox, below=of exams] (rooms) {exam\_rooms};
  \node[objbox, right=of rooms] (roomexams) {exam\_room\_exams};
  \node[clsbox, below=of rooms] (attempts) {exam\_attempts};
  \node[clsbox, left=of attempts] (submissions) {submissions};
  \node[clsbox, left=of submissions] (results) {results};
  \node[objbox, below=of users] (rubrics) {rubrics};
  \node[clsbox, below=of attempts] (classrooms) {classroom\_rooms};
  \node[objbox, left=of classrooms] (members) {classroom\_members};
  \draw[flowarr] (users) -- (questions);
  \draw[flowarr] (users) |- (exams);
  \draw[flowarr] (questions) -- (examq);
  \draw[flowarr] (exams) -- (examq);
  \draw[flowarr] (exams) -- (rooms);
  \draw[flowarr] (rooms) -- (roomexams);
  \draw[flowarr] (exams) -- (roomexams);
  \draw[flowarr] (rooms) -- (attempts);
  \draw[flowarr] (exams) -- (attempts);
  \draw[flowarr] (users) |- (attempts);
  \draw[flowarr] (attempts) -- (submissions);
  \draw[flowarr] (submissions) -- (results);
  \draw[flowarr] (exams) |- (rubrics);
  \draw[flowarr] (users) |- (classrooms);
  \draw[flowarr] (classrooms) -- (members);
  \draw[flowarr] (users) |- (members);
\end{tikzpicture}}
\caption{Sơ đồ quan hệ giữa các nhóm bảng nghiệp vụ chính.}
\label{fig:erd}
\end{figure}

\section{Thiết kế mô hình hoặc cơ sở tri thức}

Thành phần tạo câu hỏi có hai nhánh. Nhánh mô hình ngôn ngữ được gọi qua
OpenAI API khi cấu hình \code{OPENAI\_API\_KEY}; yêu cầu đầu ra gồm loại câu,
nội dung, lựa chọn, đáp án, giải thích và bằng chứng nguồn. Backend kiểm tra
định dạng, ràng buộc lựa chọn và mức độ bám nguồn trước khi nhận câu. Nhánh
dự phòng tách nội dung thành các đoạn/câu, nhận diện các mẫu kiến thức và
tạo câu hỏi theo quy tắc; nhánh này giữ hệ thống hoạt động khi dịch vụ ngoài
không sẵn sàng nhưng có phạm vi suy luận hạn chế.

Ở giai đoạn tạo đề, câu hỏi được lọc theo loại, độ khó và mức độ tư duy rồi
lưu cùng metadata đề trong cơ sở dữ liệu. Khi tạo nhiều mã đề, backend xáo
trộn thứ tự câu/lựa chọn và ánh xạ lại đáp án. Trong lúc thi, mỗi sinh viên
được gắn với một bộ câu hỏi cho lượt thi, câu trả lời và vị trí hiện tại
được lưu để hỗ trợ tiếp tục bài làm. Trắc nghiệm được tự chấm theo đáp án;
tự luận được đưa qua quy trình giảng viên chấm và ghi phản hồi.

\subsection{Thiết kế phân quyền}
Chính sách quyền được áp dụng theo vai trò và một số trường hợp theo quyền
sở hữu bản ghi. Admin quản lý dữ liệu toàn hệ thống; giảng viên thực hiện
chức năng giáo viên và tài nguyên do mình quản lý; sinh viên chỉ thao tác
với lượt thi/lớp mà mình được phép truy cập. Một token xác thực chỉ xác
nhận danh tính; backend còn phải kiểm tra quyền của tài nguyên cụ thể. RFC
7519 định nghĩa cấu trúc JWT, nhưng việc sử dụng token an toàn phụ thuộc
vào thuật toán, thời hạn, kiểm tra claim, bảo vệ khóa và thu hồi phiên
\cite{RFC7519}. Đăng nhập Google cần xác thực luồng OAuth/OpenID Connect,
kiểm tra trạng thái/redirect và liên kết tài khoản đúng theo danh tính đã
xác minh \cite{GoogleOIDC}.

\subsection{Bảo vệ API và thành phần AI}
OWASP Top 10 và OWASP API Security Top 10 nêu các nhóm rủi ro như kiểm soát
truy cập sai, lỗi xác thực, cấu hình không an toàn và khai thác tài nguyên
\cite{OWASPTop10,OWASPAPI}. Với mô hình ngôn ngữ, OWASP LLM Top 10 nhấn mạnh
prompt injection, rò rỉ thông tin và đầu ra không được kiểm soát
\cite{OWASPLLM}. Ở hệ thống thi, chính sách thực tế cần kiểm tra quyền sở
hữu từng đề/phòng, không trả đáp án đúng cho sinh viên trong khi thi, giới
hạn kích thước/định dạng tải lên, lưu khóa API ở secret store, giảm thiểu
nội dung nhạy cảm gửi ra mô hình và xem đầu ra mô hình như dữ liệu không
tin cậy. Phân tích này là yêu cầu thiết kế và không phải kết quả kiểm thử
xâm nhập.
